\documentclass[12pt,a4paper]{article}

% -------------------- Pacotes básicos --------------------
\usepackage[utf8]{inputenc}
\usepackage[T1]{fontenc}
\usepackage[brazil]{babel}
\usepackage{geometry}
\usepackage{setspace}
\usepackage{microtype}
\usepackage{hyperref}
\usepackage{graphicx}
\usepackage{booktabs}
\usepackage{array}
\usepackage{amsmath, amssymb}
\usepackage{csquotes}
\usepackage{enumitem}
\usepackage{caption}

\geometry{margin=2.5cm}
\onehalfspacing
\hypersetup{
  colorlinks=true,
  linkcolor=black,
  citecolor=black,
  urlcolor=blue
}

% -------------------- Título --------------------
\title{Mapeamento Geológico via CNN Integrando Sensoriamento Remoto e Geoquímica: \\[2pt]
adaptação metodológica e perspectivas de aplicação em áreas vegetadas}
\author{Nome do(a) autor(a) \\ \small Bacharelado em Geologia}
\date{\today}

% -------------------- Documento --------------------
\begin{document}
\maketitle

% ======================================================
\section{Introdução (Contexto e motivação)}
O mapeamento geológico desempenha papel fundamental na compreensão da estrutura e da evolução da crosta terrestre, servindo de base para estudos de prospecção mineral, geotecnia, geodinâmica e planejamento territorial. Em regiões recobertas por vegetação densa ou de difícil acesso, essa atividade enfrenta limitações consideráveis, sobretudo devido à escassez de afloramentos e às restrições impostas pelos métodos convencionais de observação direta.

Nas últimas décadas, os avanços nas tecnologias de sensoriamento remoto e nas técnicas de aprendizado de máquina transformaram a aquisição e a interpretação de dados geológicos. O uso combinado de imagens multiespectrais e hiperespectrais com levantamentos geoquímicos possibilita uma análise integrada das propriedades físicas e químicas das rochas. Entretanto, os classificadores tradicionais, como \emph{Random Forest} e \emph{Support Vector Machine}, ainda apresentam restrições quanto à interpretação espacial dos dados, resultando no chamado ``efeito sal e pimenta'' em mapas litológicos, caracterizado por classificações pontuais e pouco coerentes espacialmente.

Diante desse contexto, o artigo de Pan, Zuo e Wang (2023) propõe uma abordagem que combina dados multiespectrais e geoquímicos em um modelo baseado em redes neurais convolucionais (CNNs). Essa metodologia busca superar as limitações dos classificadores convencionais ao incorporar a correlação espacial e espectral entre as variáveis analisadas, resultando em uma interpretação mais precisa e contínua das unidades geológicas.

O presente relatório examina a metodologia proposta no artigo e discute sua aplicabilidade em diferentes contextos de mapeamento, especialmente em áreas cobertas por vegetação, enfatizando o potencial da integração entre dados geoquímicos e de sensoriamento remoto aliados a técnicas de inteligência artificial para aprimorar a identificação e a delimitação de unidades litológicas.

% ======================================================
\section{Resumo do artigo base (Fundamentação teórica)}
\textbf{Objetivo.} Pan, Zuo e Wang (2023) investigam se a integração de dados multiespectrais (ASTER, Sentinel-2A) e geoquímicos regionais --- após fusão multiorigem --- combinada a uma CNN baseada em janelas espaciais pode aumentar a acurácia e a coerência espacial do mapeamento litológico em áreas sob cobertura vegetal.

\textbf{Dados.} O estudo emprega: (i) imagens ASTER (30 m) e Sentinel-2A (10 m); (ii) 39 variáveis geoquímicas (elementos maiores e traço). Um mapa geológico serve de referência para rotular amostras e validar resultados.

\textbf{Pré-processamento e fusão.} A etapa central é a fusão \emph{Gram--Schmidt} (GS), que transfere informação espacial de maior resolução (S2A, 10 m) para a resposta espectral ASTER (30 m), preservando assinaturas espectrais e realçando textura/linhas. Em seguida, procede-se à fusão multiorigem (decomposição por pirâmide de Laplace) entre o cubo espectral fundido e as camadas geoquímicas, gerando bandas geoquímicas de alta resolução compatíveis com a malha óptica.

\textbf{Classificador.} Utiliza-se uma CNN de 10 camadas inspirada na AlexNet, com três convoluções intercaladas por três \emph{poolings}, duas camadas totalmente conectadas e saída \emph{softmax}. A entrada são \emph{patches} $21 \times 21$ empilhando bandas espectrais e variáveis geoquímicas, permitindo que o modelo aprenda contexto espacial (vizinhança) além da assinatura espectral.

\textbf{Resultados.} A acurácia global ($\sim 83\%$) supera a de um \emph{Random Forest} ($\sim 78\%$) treinado com os mesmos insumos. A CNN reduz o efeito ``sal e pimenta'' e melhora a delimitação de unidades pequenas/heterogêneas. Os ganhos decorrem da captação de padrões espaciais e da resolução efetiva maior das camadas geoquímicas fundidas.

\textbf{Limitações.} (i) Defasagem temporal entre ASTER (2003) e S2A (2020); (ii) densidade irregular de amostras geoquímicas; (iii) semelhanças espectrais entre litologias adjacentes. Ainda assim, demonstra-se a viabilidade e vantagem da abordagem CNN + fusão multiorigem para ambientes vegetados.

\textbf{Pontos de interesse para a adaptação.} (1) Tamanho da janela espacial como hiperparâmetro crítico (controla a escala geológica capturada); (2) fusão de dados para ampliar poder discriminante quando há pobreza de afloramentos; (3) arquiteturas CNN como alternativa a classificadores por-pixel, diminuindo ruído e melhorando fronteiras.

\begin{table}[htbp]
\centering
\caption{Síntese comparativa entre o artigo base e a proposta de adaptação.}
\label{tab:comparativa}
\renewcommand{\arraystretch}{1.25}
\begin{tabular}{>{\raggedright\arraybackslash}p{3.2cm} >{\raggedright\arraybackslash}p{6.1cm} >{\raggedright\arraybackslash}p{6.1cm}}
\toprule
\textbf{Aspecto} & \textbf{Artigo original} & \textbf{Proposta do grupo} \\
\midrule
Problema estudado & Mapeamento litológico em área vegetada (Jilinbaolige, China). & Mapeamento litológico em área genérica com cobertura vegetal densa (escala regional). \\
Tipo de dados & ASTER (30 m) + Sentinel-2A (10 m); 39 elementos geoquímicos; mapa geológico de referência. & Sentinel-2 L2A + ASTER; geoquímica regional (maiores e traço); opcional: DEM, magnetometria, gamaespectrometria. \\
Método (adaptação?) & Fusão GS + fusão multiorigem (pirâmide de Laplace) + CNN (patch-based $21\times 21$). & Repetir o pipeline, ajustando tamanho de \emph{patch}, bandas e hiperparâmetros; considerar UNet/ResNet-34 se houver rótulos poligonais densos. \\
Limitações & Defasagem temporal, densidade geoquímica irregular, semelhança espectral. & Geoquímica rala, heterogeneidade radiométrica, classes desbalanceadas; custo computacional. \\
Resultados obtidos / esperados & Acurácia $\approx 83\%$; redução ``sal e pimenta''; melhor identificação de unidades pequenas. & Acurácia $\geq 80\%$, kappa e F1 por classe elevadas; fronteiras mais contínuas; incerteza quantificada. \\
\bottomrule
\end{tabular}
\end{table}

% ======================================================
\section{Proposta de aplicação/adaptação}
\textbf{Novo problema/área.} Aplicar o método em região fortemente vegetada (escala 1:100\,000--1:250\,000), visando discriminar litologias em contexto de afloramento escasso e heterogeneidade estrutural.

\textbf{Dados pretendidos.}
\begin{itemize}[leftmargin=1.2cm]
  \item \textit{Óptico:} Sentinel-2 L2A (10--20 m; bandas VNIR/SWIR; índices espectrais derivados).
  \item \textit{Óptico complementar:} ASTER (VNIR/SWIR/TIR; 15/30/90 m).
  \item \textit{Geoquímica regional:} elementos maiores e traço (ICP-MS/FRX) em malha regional; interpolação se necessário.
  \item \textit{Base geológica:} mapa de referência (polígonos litológicos) para amostragem supervisionada e validação.
  \item \textit{Opcionais (recomendados):} DEM (derivados geomorfológicos), magnetometria e gamaespectrometria.
\end{itemize}

\textbf{Semelhanças e diferenças versus o artigo.} Semelhanças: tipo de problema (área vegetada), dupla fonte (espectral + geoquímica), uso de CNN com janelas. Diferenças: disponibilidade/qualidade de geoquímica; possíveis novas co-variáveis (DEM, magnetometria); sistemas de referência e escala.

\textbf{Aplicabilidade direta do método.} O \emph{pipeline} é aplicável com ajustes:
\begin{enumerate}[leftmargin=1.2cm]
  \item Pré-processamento (correção atmosférica, mosaico, registro, mascaramento de nuvens).
  \item Fusão GS (S2A$\rightarrow$ASTER) com preservação espectral.
  \item Fusão multiorigem (pirâmide de Laplace) para refinar geoquímica à malha óptica.
  \item Engenharia de atributos (índices espectrais; texturas GLCM; derivadas do DEM; filtros de lineamentos).
  \item Amostragem supervisionada (\emph{patches} centrados em polígonos bem mapeados; controle de desbalanceamento).
  \item Treino da CNN (tamanho de \emph{patch} otimizado; \emph{data augmentation}; \emph{class weights}/focal loss).
  \item Inferência (deslizante/FCN) e pós-processamento (suavização morfológica/CRF).
  \item Validação (hold-out e validação cruzada espacial; métricas por classe; matriz de confusão; PR-curves para classes raras).
  \item Quantificação de incerteza (MC-\emph{Dropout}/\emph{ensembles}) para guiar trabalho de campo.
\end{enumerate}

\textbf{Ajustes necessários.} Hiperparâmetros: tamanho de \emph{patch} (e.g., 15--31 px), \emph{batch size}, \emph{learning rate} e épocas; calibrar a escala à geologia local. Arquitetura: manter AlexNet-like para comparabilidade ou testar UNet/ResNet-34 se houver rótulos contínuos (segmentação semântica). Geoquímica rala: interpolação (IDW/krigagem) com máscara de suporte; avaliar incerteza de interpolação; ponderar por distância à amostra. Integração geofísica: empilhar bandas mag/gama normalizadas; padronizar por \emph{z-score} após imputação simples (mediana). Desbalanceamento: \emph{focal loss} ou \emph{class weights}; amostragem estratificada em polígonos menores.

\textbf{Avaliação do sucesso.} Métricas principais: acurácia global, kappa, F1 por classe (macro/micro), e IoU (se segmentação). Estratégia: validação cruzada espacial (k-fold por blocos) para reduzir \emph{leakage} espacial. Comparativos: RF/SVM por-pixel (baseline), CNN \emph{patch}-based e (havendo rótulo denso) UNet/ResNet-Seg. Verificação de campo direcionada por incerteza (amostrar fronteiras com alta entropia/baixa confiança).

\textbf{Recursos necessários.} Computação: GPU única ($\geq$ 8--12 GB) suficiente para \emph{patches} $21\times 21$; se UNet ($512^2$), ideal $\geq$16 GB; CPU para pré-processamento GDAL. Armazenamento: $\sim$50--200 GB (mosaicos, cubos fundidos, \emph{tiles} de treino). Tempo: pré-processamento (horas) e treino (horas a $\sim$1 dia, conforme arquitetura/dados).

\textbf{Fluxo (diagrama textual).}
\[
\text{Aquisição \& Pré-processamento} \;\rightarrow\; \text{Fusão GS} \;\rightarrow\; \text{Fusão multiorigem (geoquímica)}
\]
\[
\downarrow \hspace{1cm} \text{Engenharia de atributos (índices, DEM, texturas, geofísica)} \hspace{1cm} \downarrow
\]
\[
\text{Amostragem supervisionada (\emph{patches})} \;\rightarrow\; \text{Treino CNN} \;\rightarrow\; \text{Inferência} \;\rightarrow\; \text{Pós-processamento} \;\rightarrow\; \text{Validação \& Incerteza}
\]

% ======================================================
\section{Desafios e limitações}
\textbf{Disponibilidade e qualidade de dados.} Geoquímica pode ser esparsa e temporalmente não coeva com o óptico; ASTER/S2A têm datas distintas. \textit{Mitigação:} interpolação com incerteza, máscaras de confiança e priorização de bandas/épocas contemporâneas; estratificar por estação seca/úmida.

\textbf{Complexidade computacional.} Treino de CNN com muitos canais aumenta memória e tempo. \textit{Mitigação:} seleção de bandas (importância/permutação), PCA para compactação e modelos mais leves (ResNet-34 ``slim'').

\textbf{Semelhança espectral e mistura.} Litologias próximas (e.g., arenitos vs.\ riolitos alterados) geram confusão de classes. \textit{Mitigação:} índices diagnósticos (ferro/argilas), TIR/ASTER, texturas, mag/gama e janelas maiores (sem perder detalhe).

\textbf{Rotulagem imperfeita.} Mapas de referência têm incerteza cartográfica; fronteiras difusas. \textit{Mitigação:} \textit{buffer} de fronteiras na amostragem, pós-suavização (CRF/morfologia), aprendizado ruidoso (\emph{label smoothing}).

\textbf{Generalização espacial.} Risco de \emph{overfitting} quando treino e teste são vizinhos. \textit{Mitigação:} validação cruzada espacial, \emph{hold-out} por blocos e avaliação por distância.

% ======================================================
\section{Resultados esperados e perspectivas}
\textbf{Resultados esperados.}
\begin{itemize}[leftmargin=1.2cm]
  \item Acurácia $\geq 80\%$ com F1 macro superior ao \emph{baseline} por-pixel; kappa elevado.
  \item Redução do ``sal e pimenta'' e fronteiras mais contínuas.
  \item Mapas de incerteza úteis para planejamento de campo.
  \item Melhor detecção de unidades pequenas/heterogêneas.
\end{itemize}

\textbf{Contribuições.} Protocolo reprodutível de fusão e CNN para áreas vegetadas; integração efetiva de geoquímica e óptico (e, quando disponível, geofísica).

\textbf{Perspectivas.} Piloto em subárea com validação de campo; extensão para arquiteturas de segmentação (UNet/DeepLab) quando houver rótulos densos; aprendizado auto-supervisionado/pré-treino multimodal; integração de \emph{Vision Transformers} para contexto de longo alcance com dados mais abundantes.

% ======================================================
\section{Conclusão}
A adaptação do método de Pan et al.\ mostra-se viável para mapeamento litológico em ambientes vegetados, desde que se ajuste a escala do \emph{patch}, se mitigue desbalanceamento e se modele a incerteza decorrente da geoquímica esparsa e das defasagens temporais. A fusão multiorigem aliada a CNNs permite ganhos mensuráveis de acurácia e melhor continuidade espacial, posicionando o fluxo proposto como alternativa superior aos classificadores por-pixel em cenários de afloramento limitado.

% ======================================================
\section{Referências}
\begin{thebibliography}{9}

\bibitem{Pan2023}
PAN, T.; ZUO, R.; WANG, Z.
Geological Mapping via Convolutional Neural Network Based on Remote Sensing and Geochemical Survey Data in Vegetation Coverage Areas.
\textit{IEEE Journal of Selected Topics in Applied Earth Observations and Remote Sensing}, v.~16, p.~3485--3494, 2023.

\bibitem{Pohl1998}
POHL, C.; VAN GENDEREN, J. L.
Multisensor image fusion in remote sensing: concepts, methods and applications.
\textit{International Journal of Remote Sensing}, v.~19, n.~5, p.~823--854, 1998.

\bibitem{Ronneberger2015}
RONNEBERGER, O.; FISCHER, P.; BROX, T.
U-Net: Convolutional Networks for Biomedical Image Segmentation.
In: \textit{MICCAI}, 2015.

\bibitem{He2016}
HE, K.; ZHANG, X.; REN, S.; SUN, J.
Deep Residual Learning for Image Recognition.
In: \textit{CVPR}, 2016.

\end{thebibliography}

% ======================================================
\section*{Sugestões para a apresentação oral (10 min)}
\begin{itemize}[leftmargin=1.2cm]
  \item \textbf{Contexto \& problema (1 slide):} desafios em áreas vegetadas; efeito ``sal e pimenta''.
  \item \textbf{Artigo \& método (2 slides):} dados, fusões (GS + multiorigem), CNN \emph{patch}-based.
  \item \textbf{Proposta (2 slides):} dados locais, ajustes (patch, bandas, métricas), diagrama do fluxo.
  \item \textbf{Desafios \& mitigação (2 slides):} dados ralos, defasagem temporal, desbalanceamento.
  \item \textbf{Resultados esperados (1 slide):} métricas-alvo, mapas de incerteza, comparação com RF.
  \item \textbf{Conclusão (1 slide):} mensagem principal e próximos passos (piloto/validação de campo).
\end{itemize}

\end{document}
